% Curs 12 - Pagina 7 (3-I-1995).
% Fragment pentru % Curs 12 - Pagina 7 (3-I-1995).
% Fragment pentru % Curs 12 - Pagina 7 (3-I-1995).
% Fragment pentru % Curs 12 - Pagina 7 (3-I-1995).
% Fragment pentru \input{courses/pages/curs12-pagina7-ro}.
% Fara preambul si fara salt de pagina fortat.
% NOTE EDITORIALE:
% - Fragmentul continua demonstratia proprietatii universale a monoidului
%   comutativ liber, inceputa pe pagina 6.
% - Textul foarte palid din mijlocul paginii (despre grupul liber si o
%   multime A') pare a fi o incercare anterioara, stearsa sau abandonata;
%   nu a fost inclus ca text activ. Se pastreaza insa exemplul clar scris.
% - Identificarea A'={a' | a in A} este lectura contextuala a unei linii
%   palide; notatia A' apare clar in exemplul de jos.
% - Simbolul 1 din relatiile R desemneaza cuvantul vid/unitatea monoidului.
% - Literele pentru sagetile de la final sunt redate potrivit lecturii
%   aparente: zeta pentru includerea in monoidul liber si rho pentru proiectie.

\subsection*{Curs 12 - Pagina 7}

\noindent $(\mathbb{N}^{(A)},+,0)$, unde
$f:A\longrightarrow\mathbb{N}$ este de suport finit.

\medskip
\noindent $f+g\quad\Longrightarrow\quad$ monoid comutativ.

\bigskip
\noindent $x\in A$, \quad $e_x:A\longrightarrow\mathbb{N}$,
\[
e_x(y)=
\begin{cases}
1, & y=x,\\
0, & y\neq x.
\end{cases}
\]

\[
\begin{tikzcd}[column sep=huge, row sep=large]
A \arrow[r,"\eta_A"] \arrow[d,"\eta"']
  & \dfrac{A^*}{R^*} \arrow[dl,dashed,"{\exists\,g}"]\\
\mathbb{N}^{(A)} &
\end{tikzcd}
\qquad g\circ\eta_A=\eta.
\]

\noindent $\eta(x)=e_x$.
\[
\eta\text{ injectiv\u{a}}\quad\Longrightarrow\quad
\eta_A\text{ injectiv\u{a}}.
\]

\medskip
\noindent\underline{Exerci\c{t}iu}
\[
\frac{A^*}{R^*}\cong(\mathbb{N}^{(A)},+,0).
\]

\bigskip
\noindent\underline{Ex.}\quad Fie $a\in A$, 
$A'=\{a'\mid a\in A\}$.

\medskip
\noindent Consider\u{a}m monoidul liber $(A\cup A')^*$
\c{s}i rela\c{t}ia
\[
\begin{aligned}
R={}&\{(aa',1)\mid a\in A\}
  \cup\{(a'a,1)\mid a\in A\}\\
&\subseteq (A\cup A')^*\times(A\cup A')^*.
\end{aligned}
\]

\[
\mathrm{FG}[A]=\frac{(A\cup A')^*}{R^*}.
\]
\noindent\textit{Free group} (grup liber).

\medskip
\[
\begin{tikzcd}[column sep=large, row sep=large]
A \arrow[r,hook] \arrow[rrr,bend right=30,"\zeta_A"']
& A\cup A' \arrow[r,"\zeta_{A\cup A'}"]
& (A\cup A')^* \arrow[r,"\rho"]
& \dfrac{(A\cup A')^*}{R^*}
\end{tikzcd}
\]
.
% Fara preambul si fara salt de pagina fortat.
% NOTE EDITORIALE:
% - Fragmentul continua demonstratia proprietatii universale a monoidului
%   comutativ liber, inceputa pe pagina 6.
% - Textul foarte palid din mijlocul paginii (despre grupul liber si o
%   multime A') pare a fi o incercare anterioara, stearsa sau abandonata;
%   nu a fost inclus ca text activ. Se pastreaza insa exemplul clar scris.
% - Identificarea A'={a' | a in A} este lectura contextuala a unei linii
%   palide; notatia A' apare clar in exemplul de jos.
% - Simbolul 1 din relatiile R desemneaza cuvantul vid/unitatea monoidului.
% - Literele pentru sagetile de la final sunt redate potrivit lecturii
%   aparente: zeta pentru includerea in monoidul liber si rho pentru proiectie.

\subsection*{Curs 12 - Pagina 7}

\noindent $(\mathbb{N}^{(A)},+,0)$, unde
$f:A\longrightarrow\mathbb{N}$ este de suport finit.

\medskip
\noindent $f+g\quad\Longrightarrow\quad$ monoid comutativ.

\bigskip
\noindent $x\in A$, \quad $e_x:A\longrightarrow\mathbb{N}$,
\[
e_x(y)=
\begin{cases}
1, & y=x,\\
0, & y\neq x.
\end{cases}
\]

\[
\begin{tikzcd}[column sep=huge, row sep=large]
A \arrow[r,"\eta_A"] \arrow[d,"\eta"']
  & \dfrac{A^*}{R^*} \arrow[dl,dashed,"{\exists\,g}"]\\
\mathbb{N}^{(A)} &
\end{tikzcd}
\qquad g\circ\eta_A=\eta.
\]

\noindent $\eta(x)=e_x$.
\[
\eta\text{ injectiv\u{a}}\quad\Longrightarrow\quad
\eta_A\text{ injectiv\u{a}}.
\]

\medskip
\noindent\underline{Exerci\c{t}iu}
\[
\frac{A^*}{R^*}\cong(\mathbb{N}^{(A)},+,0).
\]

\bigskip
\noindent\underline{Ex.}\quad Fie $a\in A$, 
$A'=\{a'\mid a\in A\}$.

\medskip
\noindent Consider\u{a}m monoidul liber $(A\cup A')^*$
\c{s}i rela\c{t}ia
\[
\begin{aligned}
R={}&\{(aa',1)\mid a\in A\}
  \cup\{(a'a,1)\mid a\in A\}\\
&\subseteq (A\cup A')^*\times(A\cup A')^*.
\end{aligned}
\]

\[
\mathrm{FG}[A]=\frac{(A\cup A')^*}{R^*}.
\]
\noindent\textit{Free group} (grup liber).

\medskip
\[
\begin{tikzcd}[column sep=large, row sep=large]
A \arrow[r,hook] \arrow[rrr,bend right=30,"\zeta_A"']
& A\cup A' \arrow[r,"\zeta_{A\cup A'}"]
& (A\cup A')^* \arrow[r,"\rho"]
& \dfrac{(A\cup A')^*}{R^*}
\end{tikzcd}
\]
.
% Fara preambul si fara salt de pagina fortat.
% NOTE EDITORIALE:
% - Fragmentul continua demonstratia proprietatii universale a monoidului
%   comutativ liber, inceputa pe pagina 6.
% - Textul foarte palid din mijlocul paginii (despre grupul liber si o
%   multime A') pare a fi o incercare anterioara, stearsa sau abandonata;
%   nu a fost inclus ca text activ. Se pastreaza insa exemplul clar scris.
% - Identificarea A'={a' | a in A} este lectura contextuala a unei linii
%   palide; notatia A' apare clar in exemplul de jos.
% - Simbolul 1 din relatiile R desemneaza cuvantul vid/unitatea monoidului.
% - Literele pentru sagetile de la final sunt redate potrivit lecturii
%   aparente: zeta pentru includerea in monoidul liber si rho pentru proiectie.

\subsection*{Curs 12 - Pagina 7}

\noindent $(\mathbb{N}^{(A)},+,0)$, unde
$f:A\longrightarrow\mathbb{N}$ este de suport finit.

\medskip
\noindent $f+g\quad\Longrightarrow\quad$ monoid comutativ.

\bigskip
\noindent $x\in A$, \quad $e_x:A\longrightarrow\mathbb{N}$,
\[
e_x(y)=
\begin{cases}
1, & y=x,\\
0, & y\neq x.
\end{cases}
\]

\[
\begin{tikzcd}[column sep=huge, row sep=large]
A \arrow[r,"\eta_A"] \arrow[d,"\eta"']
  & \dfrac{A^*}{R^*} \arrow[dl,dashed,"{\exists\,g}"]\\
\mathbb{N}^{(A)} &
\end{tikzcd}
\qquad g\circ\eta_A=\eta.
\]

\noindent $\eta(x)=e_x$.
\[
\eta\text{ injectiv\u{a}}\quad\Longrightarrow\quad
\eta_A\text{ injectiv\u{a}}.
\]

\medskip
\noindent\underline{Exerci\c{t}iu}
\[
\frac{A^*}{R^*}\cong(\mathbb{N}^{(A)},+,0).
\]

\bigskip
\noindent\underline{Ex.}\quad Fie $a\in A$, 
$A'=\{a'\mid a\in A\}$.

\medskip
\noindent Consider\u{a}m monoidul liber $(A\cup A')^*$
\c{s}i rela\c{t}ia
\[
\begin{aligned}
R={}&\{(aa',1)\mid a\in A\}
  \cup\{(a'a,1)\mid a\in A\}\\
&\subseteq (A\cup A')^*\times(A\cup A')^*.
\end{aligned}
\]

\[
\mathrm{FG}[A]=\frac{(A\cup A')^*}{R^*}.
\]
\noindent\textit{Free group} (grup liber).

\medskip
\[
\begin{tikzcd}[column sep=large, row sep=large]
A \arrow[r,hook] \arrow[rrr,bend right=30,"\zeta_A"']
& A\cup A' \arrow[r,"\zeta_{A\cup A'}"]
& (A\cup A')^* \arrow[r,"\rho"]
& \dfrac{(A\cup A')^*}{R^*}
\end{tikzcd}
\]
.
% Fara preambul si fara salt de pagina fortat.
% NOTE EDITORIALE:
% - Fragmentul continua demonstratia proprietatii universale a monoidului
%   comutativ liber, inceputa pe pagina 6.
% - Textul foarte palid din mijlocul paginii (despre grupul liber si o
%   multime A') pare a fi o incercare anterioara, stearsa sau abandonata;
%   nu a fost inclus ca text activ. Se pastreaza insa exemplul clar scris.
% - Identificarea A'={a' | a in A} este lectura contextuala a unei linii
%   palide; notatia A' apare clar in exemplul de jos.
% - Simbolul 1 din relatiile R desemneaza cuvantul vid/unitatea monoidului.
% - Literele pentru sagetile de la final sunt redate potrivit lecturii
%   aparente: zeta pentru includerea in monoidul liber si rho pentru proiectie.

\subsection*{Curs 12 - Pagina 7}

\noindent $(\mathbb{N}^{(A)},+,0)$, unde
$f:A\longrightarrow\mathbb{N}$ este de suport finit.

\medskip
\noindent $f+g\quad\Longrightarrow\quad$ monoid comutativ.

\bigskip
\noindent $x\in A$, \quad $e_x:A\longrightarrow\mathbb{N}$,
\[
e_x(y)=
\begin{cases}
1, & y=x,\\
0, & y\neq x.
\end{cases}
\]

\[
\begin{tikzcd}[column sep=huge, row sep=large]
A \arrow[r,"\eta_A"] \arrow[d,"\eta"']
  & \dfrac{A^*}{R^*} \arrow[dl,dashed,"{\exists\,g}"]\\
\mathbb{N}^{(A)} &
\end{tikzcd}
\qquad g\circ\eta_A=\eta.
\]

\noindent $\eta(x)=e_x$.
\[
\eta\text{ injectiv\u{a}}\quad\Longrightarrow\quad
\eta_A\text{ injectiv\u{a}}.
\]

\medskip
\noindent\underline{Exerci\c{t}iu}
\[
\frac{A^*}{R^*}\cong(\mathbb{N}^{(A)},+,0).
\]

\bigskip
\noindent\underline{Ex.}\quad Fie $a\in A$, 
$A'=\{a'\mid a\in A\}$.

\medskip
\noindent Consider\u{a}m monoidul liber $(A\cup A')^*$
\c{s}i rela\c{t}ia
\[
\begin{aligned}
R={}&\{(aa',1)\mid a\in A\}
  \cup\{(a'a,1)\mid a\in A\}\\
&\subseteq (A\cup A')^*\times(A\cup A')^*.
\end{aligned}
\]

\[
\mathrm{FG}[A]=\frac{(A\cup A')^*}{R^*}.
\]
\noindent\textit{Free group} (grup liber).

\medskip
\[
\begin{tikzcd}[column sep=large, row sep=large]
A \arrow[r,hook] \arrow[rrr,bend right=30,"\zeta_A"']
& A\cup A' \arrow[r,"\zeta_{A\cup A'}"]
& (A\cup A')^* \arrow[r,"\rho"]
& \dfrac{(A\cup A')^*}{R^*}
\end{tikzcd}
\]
